%fich.tex
%%%%%%%%%%%%%%%%%%%%%%%%%%%%%%%%%%%%%%%%%%%%%%%%%%%%%%%%%%%%%%%%%%%%%%%%
% Fco. Javier Bóhorquez Ogalla
% Descripción del documento:
	
%%%%%%%%%%%%%%%%%%%%%%%%%%%%%%%%%%%%%%%%%%%%%%%%%%%%%%%%%%%%%%%%%%%%%%%%


%%%%%%%%%%%%%%%%%%%%%%%%%%%%%%%%%%%%%%%%%%%%%%%%%%%%%%%%%%%%%%%%%%%%%%%%
%Clase de documento
\documentclass[12pt, spanish]{article}
%%%%%%%%%%%%%%%%%%%%%%%%%%%%%%%%%%%%%%%%%%%%%%%%%%%%%%%%%%%%%%%%%%%%%%%%


%%%%%%%%%%%%%%%%%%%%%%%%%%%%%%%%%%%%%%%%%%%%%%%%%%%%%%%%%%%%%%%%%%%%%%%%	
%Paquetes de lenguaje:
%%%%%%%%%%%%%%%%%%%%%%%%%%%%%%%%%%%%%%%%%%%%%%%%%%%%%%%%%%%%%%%%%%%%%%%%
\usepackage[utf8]{inputenc}
\usepackage[spanish]{babel}
%\usepackage[spanish, activeacute] {babel}	
%\usepackage[spanish]{babel} 				
%\usepackage[latin1]{inputenc}
%%%%%%%%%%%%%%%%%%%%%%%%%%%%%%%%%%%%%%%%%%%%%%%%%%%%%%%%%%%%%%%%%%%%%%%%

%%%%%%%%%%%%%%%%%%%%%%%%%%%%%%%%%%%%%%%%%%%%%%%%%%%%%%%%%%%%%%%%%%%%%%%%
%Gramatic

\usepackage{/usr/share/texlive/texmf-dist/tex/latex/base/bnf}
%%%%%%%%%%%%%%%%%%%%%%%%%%%%%%%%%%%%%%%%%%%%%%%%%%%%%%%%%%%%%%%%%%%%%%%%

%%%%%%%%%%%%%%%%%%%%%%%%%%%%%%%%%%%%%%%%%%%%%%%%%%%%%%%%%%%%%%%%%%%%%%%%
%Paquetes para encabezados y pie de páginas
%%%%%%%%%%%%%%%%%%%%%%%%%%%%%%%%%%%%%%%%%%%%%%%%%%%%%%%%%%%%%%%%%%%%%%%%
\usepackage{fancyhdr}
%%%%%%%%%
% \pagestyle{fancy} %Si se cambia el estilo de página, antes de empezar 
	%el documento. Esto hace que el emcabezado y pie de página se 
	%quede separado del texto por una línea
%%%%%%%%%	
% \fancyhead{} % Límpia el texto que se está usando como encabezado
%%%%%%%%%
% \fancyhead[OPCIONES]{Encabezado}
	% OPCIONES:
		% L texto a la izquierda
		% C texto centrado
		% R texto a la derecha
		% E página par
		% O pagina impar
	%Ejemplo: 
		% fancyhead [LE] {"Doc. Latex"} 
			% Establece como emcabezado de las páginas pares "Doc Latex"
			% con alineacion a la izquierda
%%%%%%%%%			
% \fancyfoot[OPCIONES]{Píe de pag.}
	% OPCIONES:
		% L C R E O
	%Ejemplo
		% \fancyfoot[LE,RO]{\thepage} 
			% Establece como píe de pág. el número de página. con 
			% alineación a la izq en paginas pares y a la derecha en
			% las impares
%%%%%%%%%			
% \renewcommand{\headrulewidth}{0.4pt}
% \renewcommand{\footrulewidth}{0.4pt}
	% Fija el grosor de la la linea que separa el emcabezado y pie de pág
%%%%%%%%%
% Otros comandos:
%\lhead{}
%\chead{}
%\rhead{}
%\lfoot{}
%\cfoot{}
%%%%%%%%%%%%%%%%%%%%%%%%%%%%%%%%%%%%%%%%%%%%%%%%%%%%%%%%%%%%%%%%%%%%%%%%


%%%%%%%%%%%%%%%%%%%%%%%%%%%%%%%%%%%%%%%%%%%%%%%%%%%%%%%%%%%%%%%%%%%%%%%%
% Paquetes para tamaños y distancias
%%%%%%%%%%%%%%%%%%%%%%%%%%%%%%%%%%%%%%%%%%%%%%%%%%%%%%%%%%%%%%%%%%%%%%%%
\usepackage{anysize} 

%%%%%%%%%%%
% \marginsize{3cm}{2cm}{2cm}{2cm}
	% Controla los márgenes {izquierda}{derecha}{arriba}{abajo}. 
%%%%%%%%%%%%%%%%%%%%%%%%%%%%%%%%%%%%%%%%%%%%%%%%%%%%%%%%%%%%%%%%%%%%%%%%

   


%%%%%%%%%%%%%%%%%%%%%%%%%%%%%%%%%%%%%%%%%%%%%%%%%%%%%%%%%%%%%%%%%%%%%%%%
%Paquetes de carácteres especiales:
%%%%%%%%%%%%%%%%%%%%%%%%%%%%%%%%%%%%%%%%%%%%%%%%%%%%%%%%%%%%%%%%%%%%%%%%
\usepackage{dsfont}	
%Para representar conjuntos matematicos comunes: Z, N, R...
	% mathds{R}, mathds{N},...
%%%%%%%%%%%%%%%%%%%%%%%%%%%%%%%%%%%%%%%%%%%%%%%%%%%%%%%%%%%%%%%%%%%%%%%%

%%%%%%%%%%%%%%%%%%%%%%%%%%%%%%%%%%%%%%%%%%%%%%%%%%%%%%%%%%%%%%%%%%%%%%%%
%Paquetes para insertar gráficos
%%%%%%%%%%%%%%%%%%%%%%%%%%%%%%%%%%%%%%%%%%%%%%%%%%%%%%%%%%%%%%%%%%%%%%%%
\usepackage[dvips]{graphicx}
\DeclareGraphicsExtensions{.pdf,.png,.jpg} %solo para PDFLaTeX
\usepackage{graphviz}
%%%%%%%%%%%%%%%%%%%%%%%%%%%%%%%%%%%%%%%%%%%%%%%%%%%%%%%%%%%%%%%%%%%%%%%%

%%%%%%%%%%%%%%%%%%%%%%%%%%%%%%%%%%%%%%%%%%%%%%%%%%%%%%%%%%%%%%%%%%%%%%%%
%Paquetes de color:
%%%%%%%%%%%%%%%%%%%%%%%%%%%%%%%%%%%%%%%%%%%%%%%%%%%%%%%%%%%%%%%%%%%%%%%%
\usepackage{color}
%%%%%%%%%%%%%%%%%%%%%%%%%%%%%%%%%%%%%%%%%%%%%%%%%%%%%%%%%%%%%%%%%%%%%%%%
%Definición de colores:
\definecolor{gray97}{gray}{.97}
\definecolor{gray75}{gray}{.75}
\definecolor{gray45}{gray}{.45}
%%%%%%%%%%%%%%%%%%%%%%%%%%%%%%%%%%%%%%%%%%%%%%%%%%%%%%%%%%%%%%%%%%%%%%%%
\usepackage{multicol}

%%%%%%%%%%%%%%%%%%%%%%%%%%%%%%%%%%%%%%%%%%%%%%%%%%%%%%%%%%%%%%%%%%%%%%%%
%Paquete de listado de código:
%%%%%%%%%%%%%%%%%%%%%%%%%%%%%%%%%%%%%%%%%%%%%%%%%%%%%%%%%%%%%%%%%%%%%%%%
\usepackage{listings}
%%%%%%%%%%%%%%%%%%%%%%%%%%%%%%%%%%%%%%%%%%%%%%%%%%%%%%%%%%%%%%%%%%%%%%%%
%Configuración del listado:

\lstset { 
	frame=Ltb,
     framerule=0pt,
     aboveskip=0.2cm,
     framextopmargin=0.3pt,
     framexbottommargin=0.2pt,
     framexleftmargin=0.3cm,
     framesep=0pt,
     rulesep=.1pt,
     tabsize=2,
     backgroundcolor=\color{gray97},
     rulesepcolor=\color{black},
     %
     stringstyle=\ttfamily,
     showstringspaces = false,
     basicstyle=\scriptsize\ttfamily,
     commentstyle=\color{gray45},
     keywordstyle=\bfseries,
   	%
     numbers=left,
     numbersep=1pt,
     numberstyle=\tiny,
     numberfirstline = false,
     breaklines=true,
}
%%%%%%%%%%%%%%%%%%%%%%%%%%%%%%%%%%%%%%%%%%%%%%%%%%%%%%%%%%%%%%%%%%%%%%%%


%%%%%%%%%%%%%%%%%%%%%%%%%%%%%%%%%%%%%%%%%%%%%%%%%%%%%%%%%%%%%%%%%%%%%%%%
%Indexado. Índice alfabeticos
\usepackage{makeidx}
\makeindex %para habilitar indices
%%%%%%%%%%%%%%%%%%%%%%%%%%%%%%%%%%%%%%%%%%%%%%%%%%%%%%%%%%%%%%%%%%%%%%%%
%Forma de uso:
% \index{clave} %Para crear un indice de materia
% Supongase que queremos meter en el indice alfabetico referencias a 
% las paginas donde se hace referencia a la clave Producto escalar
% en tal caso en cada una de las páginas pondremos \index{Producto escalar}
% \printindex %Para imprimir el índice alfabetico

% Es necesario una doble compilación del documento. La primera genera un 
% fichero.idx el cual tendremos que procesar con la aplicación makeindex
%tras procesarlo se creará un fichero.ind que contendra el código Latex 
% que se inserta en el documento original. 
% Tras la segunda compilación del documento original se sustituira el
% comando \printindex por el contenido del fichero.ind 

%Podemos cambiar el formato de la clave:
%Ejemplo     				&    	Entrada 		&	Comentario
%\index{hola}				&         hola, 1   	&	Entrada simple
%\index{hola!Pedro}			&		 Pedro, 3 	&	Subentrada bajo ‘hola’
%\index{Juan@\textsl{Juan}}  	&		Juan, 2    	&	Entrada con diseño                        
%\index{Pepa@\textbf{Pepa}}	&		Pepa, 7    	&	Igual que antes
%\index{Loli|textbf}		&		Loli, 3    	&	No de página con diseño                     
%\index{Soraya|textit}		&		Soraya, 5  	&	Igual que antes
%%%%%%%%%%%%%%%%%%%%%%%%%%%%%%%%%%%%%%%%%%%%%%%%%%%%%%%%%%%%%%%%%%%%%%%%


%%%%%%%%%%%%%%%%%%%%%%%%%%%%%%%%%%%%%%%%%%%%%%%%%%%%%%%%%%%%%%%%%%%%%%%%


%%%%%%%%%%%%%%%%%%%%%%%%%%%%%%%%%%%%%%%%%%%%%%%%%%%%%%%%%%%%%%%%%%%%%%%%
%Estilo de página
%%%%%%%%%%%%%%%%%%%%%%%%%%%%%%%%%%%%%%%%%%%%%%%%%%%%%%%%%%%%%%%%%%%%%%%%
%\textwidth 6.75in								  %ancho de texto
%\oddsidemargin -0.2in							%margen izquierdo 
\parskip 0.2in									%espacio parrafos
%%%%%%%%%%%%%%%%%%%%%%%%%%%%%%%%%%%%%%%%%%%%%%%%%%%%%%%%%%%%%%%%%%%%%%%%
\newenvironment{changemargin}[2]{%
\begin{list}{}{%
\setlength{\topsep}{300pt}%
\setlength{\leftmargin}{#1}%
\setlength{\rightmargin}{#2}%
\setlength{\listparindent}{\parindent}%
\setlength{\itemindent}{\parindent}%
\setlength{\parsep}{\parskip}%
}%
\item[]}{\end{list}}
%%%%%%%%%%%%%%%%%%%%%%%%%%%%%%%%%%%%%%%%%%%%%%%%%%%%%%%%%%%%%%%%%%%%%%%%
\usepackage{float}
\restylefloat{table}
\usepackage{placeins}
\usepackage{framed}
\usepackage{epsf}
%%%%%%%%%%%%%%%%%%%%%%%%%%%%%%%%%%%%%%%%%%%%%%%%%%%%%%%%%%%%%%%%%%%%%%%%

\usepackage{titlesec}

\setcounter{secnumdepth}{4}

\titleformat{\paragraph}
{\normalfont\normalsize\bfseries}{\theparagraph}{1em}{}
\titlespacing*{\paragraph}
{0pt}{3.25ex plus 1ex minus .2ex}{1.5ex plus .2ex}


%%%%%%%%%%%%%%%%%%%%%%%%%%%%%%%%%%%%%%%%%%%%%%%%%%%%%%%%%%%%%%%%%%%%%%%%	
%datos del documento
\author{Gramática \\\\\ Fco. Javier Bohórquez Ogalla}						%autor
\date{}														%fecha

\title{ 
\begin{center}
\includegraphics[scale=0.5]{logo-doc.png}
\end{center} 
}
%%%%%%%%%%%%%%%%%%%%%%%%%%%%%%%%%%%%%%%%%%%%%%%%%%%%%%%%%%%%%%%%%%%%%%%%

%%%%%%%%%%%%%%%%%%%%%%%%%%%%%%%%%%%%%%%%%%%%%%%%%%%%%%%%%%%%%%%%%%%%%%%%
%Documento:
%%%%%%%%%%%%%%%%%%%%%%%%%%%%%%%%%%%%%%%%%%%%%%%%%%%%%%%%%%%%%%%%%%%%%%%%
\lstdefinestyle{nonumbers}{numbers=none}
\begin{document}

\maketitle
\pagebreak
\tableofcontents
\pagebreak
\digraph[scale=0.5]{MyGraph}{
 node [fontname = courier, shape = box, colorscheme = paired6]
  edge [fontname = courier]

  0 [label="State 0\n\l  0 $accept: . program $end\l"]
  0 -> 1 [style=dotted]
  0 -> 2 [style=solid label="NUM"]
  0 -> 3 [style=solid label="SNUM"]
  0 -> 4 [style=solid label="REXP"]
  0 -> 5 [style=solid label="STR"]
  0 -> 6 [style=solid label="NULLTOKEN"]
  0 -> 7 [style=solid label="TRUE"]
  0 -> 8 [style=solid label="FALSE"]
  0 -> 9 [style=solid label="ID"]
  0 -> 10 [style=solid label="PRINT"]
  0 -> 11 [style=solid label="INC"]
  0 -> 12 [style=solid label="DEC"]
  0 -> 13 [style=solid label="INCLUDE"]
  0 -> 14 [style=solid label="GLOBAL"]
  0 -> 15 [style=solid label="ISNULL"]
  0 -> 16 [style=solid label="FUNCTION"]
  0 -> 17 [style=solid label="PARTIALFUNC"]
  0 -> 18 [style=solid label="RET"]
  0 -> 19 [style=solid label="DECORATOR"]
  0 -> 20 [style=solid label="CONTEXTFUNC"]
  0 -> 21 [style=solid label="CLASS"]
  0 -> 22 [style=solid label="NEW"]
  0 -> 23 [style=solid label="THIS"]
  0 -> 24 [style=solid label="PARENT"]
  0 -> 25 [style=solid label="PRIVATE"]
  0 -> 26 [style=solid label="STATIC"]
  0 -> 27 [style=solid label="IF"]
  0 -> 28 [style=solid label="WITH"]
  0 -> 29 [style=solid label="DO"]
  0 -> 30 [style=solid label="WHILE"]
  0 -> 31 [style=solid label="FOR"]
  0 -> 32 [style=solid label="GOTO"]
  0 -> 33 [style=solid label="SWITCH"]
  0 -> 34 [style=solid label="TYPEOF"]
  0 -> 35 [style=solid label="SIZEOF"]
  0 -> 36 [style=solid label="DATINFO"]
  0 -> 37 [style=solid label="BREAK"]
  0 -> 38 [style=solid label="EXIT"]
  0 -> 39 [style=solid label="SIZE"]
  0 -> 40 [style=solid label="EMPTY"]
  0 -> 41 [style=solid label="CONTINUE"]
  0 -> 42 [style=solid label="TRY"]
  0 -> 43 [style=solid label="THROW"]
  0 -> 44 [style=solid label="FLOW"]
  0 -> 45 [style=solid label="EXPLODE"]
  0 -> 46 [style=solid label="IMPLODE"]
  0 -> 47 [style=solid label="SPRINTF"]
  0 -> 48 [style=solid label="FIND"]
  0 -> 49 [style=solid label="REPLACE"]
  0 -> 50 [style=solid label="SUBREPLACE"]
  0 -> 51 [style=solid label="UPPER"]
  0 -> 52 [style=solid label="LOWER"]
  0 -> 53 [style=solid label="DATE"]
  0 -> 54 [style=solid label="TIME"]
  0 -> 55 [style=solid label="SLEEP"]
  0 -> 56 [style=solid label="GETENV"]
  0 -> 57 [style=solid label="REDUCE"]
  0 -> 58 [style=solid label="FTOKEN"]
  0 -> 59 [style=solid label="FPUT"]
  0 -> 60 [style=solid label="FGET"]
  0 -> 61 [style=solid label="FSEEK"]
  0 -> 62 [style=solid label="FTELL"]
  0 -> 63 [style=solid label="FCLOSE"]
  0 -> 64 [style=solid label="FREAD"]
  0 -> 65 [style=solid label="FWRITE"]
  0 -> 66 [style=solid label="FAPP"]
  0 -> 67 [style=solid label="FORK"]
  0 -> 68 [style=solid label="WAIT"]
  0 -> 69 [style=solid label="GETPID"]
  0 -> 70 [style=solid label="GETPPID"]
  0 -> 71 [style=solid label="PROCESSTOKEN"]
  0 -> 72 [style=solid label="EXITPROCESS"]
  0 -> 73 [style=solid label="SIGNAL"]
  0 -> 74 [style=solid label="SHANDLER"]
  0 -> 75 [style=solid label="INPUT"]
  0 -> 76 [style=solid label="INPUTLINE"]
  0 -> 77 [style=solid label="MATCH"]
  0 -> 78 [style=solid label="REGEXP"]
  0 -> 79 [style=solid label="SEARCH"]
  0 -> 80 [style=solid label="EXEC"]
  0 -> 81 [style=solid label="EVAL"]
  0 -> 82 [style=solid label="ARRAYCHUNK"]
  0 -> 83 [style=solid label="LOAD_MODULE"]
  0 -> 84 [style=solid label="DEF"]
  0 -> 85 [style=solid label="CDATA"]
  0 -> 86 [style=solid label="DDATA"]
  0 -> 87 [style=solid label="SDATA"]
  0 -> 88 [style=solid label="EXIST"]
  0 -> 89 [style=solid label="CINT"]
  0 -> 90 [style=solid label="CFLOAT"]
  0 -> 91 [style=solid label="CBOOL"]
  0 -> 92 [style=solid label="CSTR"]
  0 -> 93 [style=solid label="';'"]
  0 -> 94 [style=solid label="'('"]
  0 -> 95 [style=solid label="'{'"]
  0 -> 96 [style=solid label="'$'"]
  0 -> 97 [style=solid label="'~'"]
  0 -> 98 [style=solid label="'!'"]
  0 -> 99 [style=solid label="'+'"]
  0 -> 100 [style=solid label="'-'"]
  0 -> 101 [style=solid label="'['"]
  0 -> 102 [style=dashed label="program"]
  0 -> 103 [style=dashed label="stmts"]
  0 -> 104 [style=dashed label="stmtb"]
  0 -> 105 [style=dashed label="stmt"]
  0 -> 106 [style=dashed label="label"]
  0 -> 107 [style=dashed label="define"]
  0 -> 108 [style=dashed label="if"]
  0 -> 109 [style=dashed label="with"]
  0 -> 110 [style=dashed label="while"]
  0 -> 111 [style=dashed label="dowhile"]
  0 -> 112 [style=dashed label="for"]
  0 -> 113 [style=dashed label="break"]
  0 -> 114 [style=dashed label="switch"]
  0 -> 115 [style=dashed label="iloop"]
  0 -> 116 [style=dashed label="trycatch"]
  0 -> 117 [style=dashed label="function"]
  0 -> 118 [style=dashed label="class"]
  0 -> 119 [style=dashed label="component"]
  0 -> 120 [style=dashed label="exp"]
  0 -> 121 [style=dashed label="lexp"]
  0 -> 122 [style=dashed label="lop1"]
  0 -> 123 [style=dashed label="lop2"]
  0 -> 124 [style=dashed label="lop3"]
  0 -> 125 [style=dashed label="lop4"]
  0 -> 126 [style=dashed label="aexp"]
  0 -> 127 [style=dashed label="aop1"]
  0 -> 128 [style=dashed label="aop2"]
  0 -> 129 [style=dashed label="aop3"]
  0 -> 130 [style=dashed label="sexp"]
  0 -> 131 [style=dashed label="sop1"]
  0 -> 132 [style=dashed label="sop2"]
  0 -> 133 [style=dashed label="call_identity"]
  0 -> 134 [style=dashed label="extra_operator"]
  0 -> 135 [style=dashed label="tern"]
  0 -> 136 [style=dashed label="notnull"]
  0 -> 137 [style=dashed label="unity"]
  0 -> 138 [style=dashed label="inc_dec"]
  0 -> 139 [style=dashed label="assignation"]
  0 -> 140 [style=dashed label="cexp"]
  0 -> 141 [style=dashed label="lop_unity"]
  0 -> 142 [style=dashed label="aop_unity"]
  0 -> 143 [style=dashed label="sop_unity"]
  0 -> 144 [style=dashed label="class_exp"]
  0 -> 145 [style=dashed label="arrayexp"]
  0 -> 146 [style=dashed label="date"]
  0 -> 147 [style=dashed label="file"]
  0 -> 148 [style=dashed label="process"]
  0 -> 149 [style=dashed label="environment"]
  0 -> 150 [style=dashed label="exp_generator"]
  0 -> 151 [style=dashed label="regexp_op"]
  0 -> 152 [style=dashed label="iloop_access"]
  0 -> 153 [style=dashed label="component_unity"]
  0 -> 154 [style=dashed label="identity"]
  0 -> 155 [style=dashed label="function_exp"]
  0 -> 156 [style=dashed label="function_lambda"]
  0 -> 157 [style=dashed label="function_decorator"]
  0 -> 158 [style=dashed label="function_partial"]
  0 -> 159 [style=dashed label="assig"]
  0 -> 160 [style=dashed label="parray"]
  0 -> 161 [style=dashed label="arrayop"]
  0 -> 162 [style=dashed label="namespace"]
  0 -> 163 [style=dashed label="id"]
  0 -> "0R2" [label="[$end]", style=solid]
 "0R2" [label="R2", fillcolor=3, shape=diamond, style=filled]
  
}


\end{document}
%%%%%%%%%%%%%%%%%%%%%%%%%%%%%%%%%%%%%%%%%%%%%%%%%%%%%%%%%%%%%%%%%%%%%%%%
%Fco. Javier Bohórquez Ogalla
